% GENERATED FILE. Edit canonical lessons, then run npm run generate:books.
% canonical-source-hash: fnv1a64:72a0a9c2ea0b0a5f
% canonical-lessons: TA-C373-araciyal, TA-C373-amaiccar, TA-C373-piratamar, TA-C373-alunar, TA-C373-natalumanram

\chapter{Government}
\label{ch:ta-government}
\hlchaptermodality{373}

\begin{chapteropening}
\textbf{By the end of this chapter:} \emph{I can talk about politics, ministers and parliament.}

Complete the last lesson of chapter 373: I can talk about politics, ministers and parliament.
\end{chapteropening}

\section[\texorpdfstring{\ta{அரசியல்} (araciyal)}{அரசியல் (araciyal)}]{\ta{அரசியல்} (araciyal) — politics}

\label{lesson:TA-C373-araciyal}



\begin{quote}
Before the new one: say the Tamil for to punish, then the Tamil for to comply.
\end{quote}

\subsection*{You'll want to know: \ta{அரசியல்}}
\textbf{\ta{அரசியல்}} — \emph{araciyal} — \textquotedblleft{}politics\textquotedblright{}.

\begin{quote}
\ta{என்} \ta{தாத்தாவுக்கு} \ta{அரசியல்} \ta{பற்றிப்} \ta{பேசப்} \ta{பிடிக்கும்}.
\end{quote}

\begin{quote}
\emph{eṉ tāttāvukku araciyal paṟṟip pēcap piṭikkum.}
\end{quote}

\begin{quote}
My grandfather likes to talk about politics.
\end{quote}

\ta{அரசு}, government + -\ta{இயல்}, a field: \ta{அரசியல்வாதி}, a politician.

\begin{cousinweb}
\textbf{\ta{அரசியல்}} (\emph{araciyal}) is built from \emph{aracu}, \textquotedblleft{}government\textquotedblright{}, which DEDR 201 traces to Sanskrit \emph{rājan}, \textquotedblleft{}king\textquotedblright{} (CDIAL 10679), and native \emph{iyal}, \textquotedblleft{}nature, field of study\textquotedblright{} (DEDR 471).
\end{cousinweb}

\begin{grammarlens}[title={What you've built}]
One more word for giving reasons and opinions.
\end{grammarlens}

\subsection*{Your turn}
\begin{itemize}
\raggedright
  \item \emph{Say it:} \emph{araciyal}
  \item \emph{Say it:} \emph{araciyal}, once more
  \item \emph{Say it:} \emph{eṉ tāttāvukku araciyal paṟṟip pēcap piṭikkum.} — My grandfather likes to talk about politics.
  \item \emph{Recall:} say the Tamil for to punish, then the Tamil for to comply, then say \emph{araciyal} again
\end{itemize}

\begin{tcolorbox}[breakable,colback=teal!4,colframe=teal!35!black,title={Before you move on}]
What does \ta{அரசியல்} mean? (\textquotedblleft{}Politics\textquotedblright{}.) Say it once more.
\end{tcolorbox}
\section[\texorpdfstring{\ta{அமைச்சர்} (amaiccar)}{அமைச்சர் (amaiccar)}]{\ta{அமைச்சர்} (amaiccar) — a minister}

\label{lesson:TA-C373-amaiccar}



\begin{quote}
Before the new one: say the Tamil for to comply, then the Tamil for politics.
\end{quote}

\subsection*{You'll want to know: \ta{அமைச்சர்}}
\textbf{\ta{அமைச்சர்}} — \emph{amaiccar} — \textquotedblleft{}a minister\textquotedblright{}.

\begin{quote}
\ta{கல்வி} \ta{அமைச்சர்} \ta{எங்கள்} \ta{பள்ளிக்கு} \ta{வந்தார்}.
\end{quote}

\begin{quote}
\emph{kalvi amaiccar eṅkaḷ paḷḷikku vantār.}
\end{quote}

\begin{quote}
The education minister came to our school.
\end{quote}

\ta{அமைச்சகம்}, a ministry.

\begin{cousinweb}
\textbf{\ta{அமைச்சர்}} (\emph{amaiccar}) is usually traced to native \emph{amai}, \textquotedblleft{}to be fitted, joined, settled\textquotedblright{} (DEDR 162), through \emph{amaiccu}, \textquotedblleft{}counsel\textquotedblright{}: one who sets the affairs of state in order.
\end{cousinweb}

\begin{grammarlens}[title={What you've built}]
One more word for giving reasons and opinions.
\end{grammarlens}

\subsection*{Your turn}
\begin{itemize}
\raggedright
  \item \emph{Say it:} \emph{amaiccar}
  \item \emph{Say it:} \emph{amaiccar}, once more
  \item \emph{Say it:} \emph{kalvi amaiccar eṅkaḷ paḷḷikku vantār.} — The education minister came to our school.
  \item \emph{Recall:} say the Tamil for to comply, then the Tamil for politics, then say \emph{amaiccar} again
\end{itemize}

\begin{tcolorbox}[breakable,colback=teal!4,colframe=teal!35!black,title={Before you move on}]
What does \ta{அமைச்சர்} mean? (\textquotedblleft{}A minister\textquotedblright{}.) Say it once more.
\end{tcolorbox}
\section[\texorpdfstring{\ta{பிரதமர்} (piratamar)}{பிரதமர் (piratamar)}]{\ta{பிரதமர்} (piratamar) — a prime minister}

\label{lesson:TA-C373-piratamar}



\begin{quote}
Before the new one: say the Tamil for politics, then the Tamil for a minister.
\end{quote}

\subsection*{You'll want to know: \ta{பிரதமர்}}
\textbf{\ta{பிரதமர்}} — \emph{piratamar} — \textquotedblleft{}a prime minister\textquotedblright{}.

\begin{quote}
\ta{பிரதமர்} \ta{இன்று} \ta{நாட்டு} \ta{மக்களிடம்} \ta{பேசினார்}.
\end{quote}

\begin{quote}
\emph{piratamar iṉṟu nāṭṭu makkaḷiṭam pēciṉār.}
\end{quote}

\begin{quote}
The prime minister spoke to the nation today.
\end{quote}

\ta{பிரதம}, chief, first + -\ta{அர்}, the person.

\begin{cousinweb}
\textbf{\ta{பிரதமர்}} (\emph{piratamar}) is built on Sanskrit \emph{prathama}, \textquotedblleft{}first\textquotedblright{} (CDIAL 8648), with the personal ending -\ta{அர்}: the first minister.
\end{cousinweb}

\begin{grammarlens}[title={What you've built}]
One more word for giving reasons and opinions.
\end{grammarlens}

\subsection*{Your turn}
\begin{itemize}
\raggedright
  \item \emph{Say it:} \emph{piratamar}
  \item \emph{Say it:} \emph{piratamar}, once more
  \item \emph{Say it:} \emph{piratamar iṉṟu nāṭṭu makkaḷiṭam pēciṉār.} — The prime minister spoke to the nation today.
  \item \emph{Recall:} say the Tamil for politics, then the Tamil for a minister, then say \emph{piratamar} again
\end{itemize}

\begin{tcolorbox}[breakable,colback=teal!4,colframe=teal!35!black,title={Before you move on}]
What does \ta{பிரதமர்} mean? (\textquotedblleft{}A prime minister\textquotedblright{}.) Say it once more.
\end{tcolorbox}
\section[\texorpdfstring{\ta{ஆளுநர்} (āḷunar)}{ஆளுநர் (āḷunar)}]{\ta{ஆளுநர்} (āḷunar) — a governor}

\label{lesson:TA-C373-alunar}



\begin{quote}
Before the new one: say the Tamil for a minister, then the Tamil for a prime minister.
\end{quote}

\subsection*{You'll want to know: \ta{ஆளுநர்}}
\textbf{\ta{ஆளுநர்}} — \emph{āḷunar} — \textquotedblleft{}a governor\textquotedblright{}.

\begin{quote}
\ta{ஆளுநர்} \ta{இன்று} \ta{விழாவில்} \ta{கலந்துகொண்டார்}.
\end{quote}

\begin{quote}
\emph{āḷunar iṉṟu viḻāvil kalantukoṇṭār.}
\end{quote}

\begin{quote}
The governor attended the ceremony today.
\end{quote}

\ta{ஆள்}, to rule + -\ta{நர்}, the one who.

\begin{cousinweb}
\textbf{\ta{ஆளுநர்}} (\emph{āḷunar}) is native: \emph{āḷ}, \textquotedblleft{}to rule, govern\textquotedblright{} (DEDR 5157), with the agent ending -\ta{நர்}: \textquotedblleft{}one who governs\textquotedblright{}.
\end{cousinweb}

\begin{grammarlens}[title={What you've built}]
One more word for giving reasons and opinions.
\end{grammarlens}

\subsection*{Your turn}
\begin{itemize}
\raggedright
  \item \emph{Say it:} \emph{āḷunar}
  \item \emph{Say it:} \emph{āḷunar}, once more
  \item \emph{Say it:} \emph{āḷunar iṉṟu viḻāvil kalantukoṇṭār.} — The governor attended the ceremony today.
  \item \emph{Recall:} say the Tamil for a minister, then the Tamil for a prime minister, then say \emph{āḷunar} again
\end{itemize}

\begin{tcolorbox}[breakable,colback=teal!4,colframe=teal!35!black,title={Before you move on}]
What does \ta{ஆளுநர்} mean? (\textquotedblleft{}A governor\textquotedblright{}.) Say it once more.
\end{tcolorbox}
\section[\texorpdfstring{\ta{நாடாளுமன்றம்} (nāṭāḷumaṉṟam)}{நாடாளுமன்றம் (nāṭāḷumaṉṟam)}]{\ta{நாடாளுமன்றம்} (nāṭāḷumaṉṟam) — a parliament}

\label{lesson:TA-C373-natalumanram}



\begin{quote}
Before the new one: say the Tamil for a prime minister, then the Tamil for a governor.
\end{quote}

\subsection*{You'll want to know: \ta{நாடாளுமன்றம்}}
\textbf{\ta{நாடாளுமன்றம்}} — \emph{nāṭāḷumaṉṟam} — \textquotedblleft{}a parliament\textquotedblright{}.

\begin{quote}
\ta{நாடாளுமன்றத்தில்} \ta{புதிய} \ta{சட்டம்} \ta{பற்றி} \ta{விவாதித்தார்கள்}.
\end{quote}

\begin{quote}
\emph{nāṭāḷumaṉṟattil putiya caṭṭam paṟṟi vivātittārkaḷ.}
\end{quote}

\begin{quote}
They debated the new law in parliament.
\end{quote}

\ta{நாடு}, country, + \ta{ஆளும்}, ruling, + a word for a hall.

\begin{cousinweb}
\textbf{\ta{நாடாளுமன்றம்}} (\emph{nāṭāḷumaṉṟam}) is a modern native coinage: \emph{nāṭu}, \textquotedblleft{}country\textquotedblright{} (DEDR 3638), \emph{āḷum}, \textquotedblleft{}that rules\textquotedblright{}, from \emph{āḷ} (DEDR 5157), and \emph{maṉṟam}, \textquotedblleft{}hall, assembly\textquotedblright{} (DEDR 4777).
\end{cousinweb}

\begin{grammarlens}[title={What you've built}]
One more word for giving reasons and opinions.
\end{grammarlens}

\subsection*{Your turn}
\begin{itemize}
\raggedright
  \item \emph{Say it:} \emph{nāṭāḷumaṉṟam}
  \item \emph{Say it:} \emph{nāṭāḷumaṉṟam}, once more
  \item \emph{Say it:} \emph{nāṭāḷumaṉṟattil putiya caṭṭam paṟṟi vivātittārkaḷ.} — They debated the new law in parliament.
  \item \emph{Recall:} say the Tamil for a prime minister, then the Tamil for a governor, then say \emph{nāṭāḷumaṉṟam} again
\end{itemize}

\begin{tcolorbox}[breakable,colback=teal!4,colframe=teal!35!black,title={Before you move on}]
What does \ta{நாடாளுமன்றம்} mean? (\textquotedblleft{}A parliament\textquotedblright{}.) Say it once more.
\end{tcolorbox}
